% 28 keys of a document and how a diagram reaches them: \tikzset against \tikzcdset, paths on keys, styles with a value,
% .default, .initial, key names over lines, keys that do nothing, overlay and scale.
\documentclass{article}
\pdfcompresslevel=0 \pdfobjcompresslevel=0
\usepackage{tikz-cd}
\usetikzlibrary{decorations.pathmorphing}
\tikzset{
  tdash/.style={dashed},
  tthick/.style={thick,dotted},
  tsimple/.style={->},
  tcd/.style={/tikz/commutative diagrams/two heads},
  tdouble/.style={double},
  tval/.style={line width=#1},
  tdef/.style={draw=#1},
  tdef/.default=red,
  symbol/.style={"#1" description, allow upside down, auto=false, draw=none, sloped},
  sym/.style={symbol={x}},
  tcell/.style={inner sep=#1},
  symbol/.style={"#1" description, allow upside down, auto=false, draw=none, sloped},
  sym/.style={symbol={x}},
  tcell/.style={inner sep=#1},
  tapp/.style={thick},
  tapp/.append style={red},
  unused/.code={\relax},
  unused2/.style 2 args={#1,#2},
  own/.cd, value/.initial=3, other/.initial=x,
}
\tikzcdset{
  cdash/.style={dashed},
  chead/.style={two heads,shift left=2pt},
  cval/.style={shift left=#1,"#1"},
  ccell/.style={inner sep=#1},
  clabel/.style={near end,red},
  cbig/.style={row sep=#1,column sep=#1},
  row sep/roomy/.initial=3cm,
  column sep/roomy/.initial={2cm,between origins},
  every label/.append style={},
}
\begin{document}
% in /tikz, tikz-cd's dashed is out of reach
\begin{tikzcd}A \arrow[r,tdash] & B\end{tikzcd}
\begin{tikzcd}A \arrow[r,cdash] & B\end{tikzcd}
\begin{tikzcd}A \arrow[r,tthick] & B\end{tikzcd}
\begin{tikzcd}A \arrow[r,tsimple] & B\end{tikzcd}
\begin{tikzcd}A \arrow[r,tcd] & B\end{tikzcd}
\begin{tikzcd}A \arrow[r,tdouble] & B\end{tikzcd}
\begin{tikzcd}A \arrow[r,chead] & B\end{tikzcd}
% a path on a key
\begin{tikzcd}A \arrow[r,/tikz/dashed] & B\end{tikzcd}
\begin{tikzcd}A \arrow[r,/tikz/commutative diagrams/dashed] & B\end{tikzcd}
\begin{tikzcd}A \arrow[r,/tikz/thick,/tikz/commutative diagrams/two heads] & B\end{tikzcd}
\begin{tikzcd}A \arrow[r,"f"{/tikz/near end}] & B\end{tikzcd}
\begin{tikzcd}|[/tikz/inner sep=1pt]|A \arrow[r] & B\end{tikzcd}
% values
\begin{tikzcd}A \arrow[r,tval=1pt] & B\end{tikzcd}
\begin{tikzcd}A \arrow[r,tval=2pt] \arrow[d,tval=1pt] & B\\C\end{tikzcd}
\begin{tikzcd}A \arrow[r,tdef] & B\end{tikzcd}
\begin{tikzcd}A \arrow[r,tdef=blue] & B\end{tikzcd}
\begin{tikzcd}A \arrow[r,tapp] & B\end{tikzcd}
\begin{tikzcd}A \arrow[r,cval=3] & B\end{tikzcd}
\begin{tikzcd}A \arrow[r,cval=1pt] \arrow[d,cval=2] & B\\C\end{tikzcd}
\begin{tikzcd}|[tcell=1pt]|A \arrow[r] & B\end{tikzcd}
\begin{tikzcd}[cells={ccell=1pt}]A \arrow[r] & B\end{tikzcd}
\begin{tikzcd}[every cell/.append style={ccell=1pt}]A \arrow[r] & B\end{tikzcd}
\begin{tikzcd}[cells={tcell=1pt}]A \arrow[r] & B\end{tikzcd}
\begin{tikzcd}A \arrow[r,"f"{clabel}] & B\end{tikzcd}
\begin{tikzcd}[cbig=4em]A \arrow[r] \arrow[d] & B\\C\end{tikzcd}
% names of separations
\begin{tikzcd}[row sep=roomy,column sep=roomy]A \arrow[r] \arrow[d] & B\\C\end{tikzcd}
% key names over lines, and with more spaces
\begin{tikzcd}A \arrow[r,bend
   left] \arrow[d,shift    right] & B\\C\end{tikzcd}
\begin{tikzcd}A \arrow[r,"f"{near
   end}] & B\end{tikzcd}
% what does nothing
\begin{tikzcd}A \arrow[r,left,"f"] & B\end{tikzcd}
\begin{tikzcd}A \arrow[r,"f",right] & B\end{tikzcd}
\begin{tikzcd}A \arrow[r,"f",above] & B\end{tikzcd}
\begin{tikzcd}A \arrow[r,"f",below] & B\end{tikzcd}
\begin{tikzcd}[arrows=left]A \arrow[r,"f"] & B\end{tikzcd}
\begin{tikzcd}A \arrow[r,"f",auto] & B\end{tikzcd}
\begin{tikzcd}A \arrow[r,"f",auto=false] & B\end{tikzcd}
\begin{tikzcd}A \arrow[r,sym,red] & B\end{tikzcd}
\begin{tikzcd}A \arrow[dr,sym] \\ & B\end{tikzcd}
\begin{tikzcd}A \arrow[r,"f",auto=false] & B\end{tikzcd}
\begin{tikzcd}A \arrow[r,sym,red] & B\end{tikzcd}
\begin{tikzcd}A \arrow[dr,sym] \\ & B\end{tikzcd}
\begin{tikzcd}[every node/.style={red}]A \arrow[r,"f"] & B\end{tikzcd}
\begin{tikzcd}[scale=1]A \arrow[r,shift left,"f"] & B\end{tikzcd}
% overlay: nothing of it counts for the box
\begin{tikzcd}A \arrow[r,overlay] & B\end{tikzcd}
\begin{tikzcd}A \arrow[r,overlay,"f"] & B\end{tikzcd}
\begin{tikzcd}A \arrow[rr,overlay,bend left=80] & & B\end{tikzcd}
\begin{tikzcd}A \arrow[r,"f"{overlay}] & B\end{tikzcd}
% a label's own scale: its shape too, and what follows in its keys
\begin{tikzcd}A \arrow[r,"f"{scale=2}] & B\end{tikzcd}
\begin{tikzcd}A \arrow[r,"f"{scale=2,xshift=1cm}] & B\end{tikzcd}
\begin{tikzcd}A \arrow[r,"f"{xshift=1cm,scale=2}] & B\end{tikzcd}
\begin{tikzcd}A \arrow[r,"f"{scale=2,above=3pt}] & B\end{tikzcd}
\begin{tikzcd}A \arrow[r,"f"{scale=1.5,description}] & B\end{tikzcd}
% colours, as they reach the operators
\begin{tikzcd}A \arrow[r,red,"f"] & B\end{tikzcd}
\begin{tikzcd}A \arrow[r,color=red,"f"] & B\end{tikzcd}
\begin{tikzcd}A \arrow[r,draw=green,"f"] & B\end{tikzcd}
\begin{tikzcd}A \arrow[r,red,blue,"f"] & B\end{tikzcd}
\begin{tikzcd}A \arrow[r,"f"{red}] & B\end{tikzcd}
\begin{tikzcd}A \arrow[r,"f"{color=red}] & B\end{tikzcd}
\begin{tikzcd}A \arrow[r,"f"{text=blue}] & B\end{tikzcd}
\begin{tikzcd}A \arrow[r,text=blue,"f"] & B\end{tikzcd}
\begin{tikzcd}A \arrow[r,red,"f"{description}] & B\end{tikzcd}
\begin{tikzcd}A \arrow[r,"f"{fill=yellow}] & B\end{tikzcd}
\begin{tikzcd}|[red]|A \arrow[r] & |[draw]|B\end{tikzcd}
\begin{tikzcd}|[draw=red,rounded corners]|A \arrow[r] & |[color=blue]|B\end{tikzcd}
\begin{tikzcd}[cells={nodes={red}}]A \arrow[r] & |[blue]|B\end{tikzcd}
\begin{tikzcd}A \arrow[r,solid,dashed,solid] & B\end{tikzcd}
\end{document}
